% 2de graad, Eerstegraadsvergelijkingen (reeks Vergelijkingen, met Stelsels).
% De Hooch: een vrouw met een weegschaal. Wat je links doet, doe je rechts.
% Motief: een balans in evenwicht, 2x + 1 = 7, dus x = 3.
\def\omslagbestand{omslag-goudweegster.jpg}
\def\omslagbron{Pieter de Hooch, \textit{De goudweegster}, ca. 1664 (publiek domein)}
\def\omslagrand{10.7 + 0.3*sin(\x*360/14)}
\def\omslagmotief{%
  \randlijn
  \draw[white,line width=1.6pt] (16.3,4.7) -- (17.7,4.7);
  \draw[white,line width=1.6pt] (17,4.7) -- (17,7.7);
  \draw[white,line width=1.6pt] (14.8,7.7) -- (19.2,7.7);
  \fill[omslagstip] (17,7.7) circle (1.3mm);
  \foreach \m in {15.2,18.8}{%
    \draw[white,line width=.8pt] (\m,7.7) -- (\m-.75,6.3);
    \draw[white,line width=.8pt] (\m,7.7) -- (\m+.75,6.3);
    \draw[white,line width=1.6pt] (\m-.85,6.3) -- (\m+.85,6.3);
  }
  \node[fill=fel,text=inkt,rounded corners=1.5pt,inner sep=2.5pt,anchor=south] at (15.2,6.35) {$2x+1$};
  \node[fill=fel,text=inkt,rounded corners=1.5pt,inner sep=2.5pt,anchor=south] at (18.8,6.35) {$7$};
  \node[text=fel,font=\large,anchor=west] at (17.9,5.0) {$x=3$};
}
